\section{Series de Fourier}
\label{sec:fourier}

Supongamos que $f$ es continua por pedazos en $x\in[a,b]$ y suave, es decir,
que $\dv{}{x}f(x)$ es continua por pedazos en $(a,b)$, y pidamos adem\'as que
existan $\eval{\dv{}{x}f(x)}_{a^{+}}$ y $\eval{\dv{}{x}f(x)}_{b^{-}}$. De ese
universo de funciones nos interesan las peri\'odicas de periodo $2\pi$.

Si $f$ admite una expansi\'on en alguna base $\set{g_k}$ como en
\eqref{eq:expansion_general}, $f(x)=\sum_k C_kg_k(x)$, la elecci\'on
\begin{equation}
  \label{eq:base_fourier}
  g_k(x)=\qty{e^{ix}}^{k}=e^{ikx}
\end{equation}
produce la \textbf{serie de Fourier}
\begin{equation}
  \label{eq:serie_fourier}
  f(x)=\sum_{k=-\infty}^{\infty}C_ke^{ikx} .
\end{equation}

\begin{teorema}[Convergencia]
  \label{teo:convergencia_fourier}
  Si $f$ es peri\'odica de periodo $2\pi$, continua y suave por pedazos,
  entonces su serie de Fourier converge absoluta y uniformemente a $f$.
\end{teorema}

\subsection{Par\'entesis: producto interno de funciones}
\label{subsec:producto_interno_funciones}

Para vectores, el producto interno es
$\braket{\vb{u}}{\vb{v}}=\vb{u}^\dg\vb{v}$. Para funciones se define de
manera an\'aloga sustituyendo la suma sobre componentes por una integral,
\begin{equation}
  \label{eq:pi_funciones}
  \braket{g}{f}=\frac{1}{N}\int_{a}^{b}g^{*}(x)f(x)\dd x,
\end{equation}
donde $N$ es una constante de normalizaci\'on. Apliqu\'emoslo a la base
\eqref{eq:base_fourier}:
\begin{equation}
  \label{eq:pi_exponenciales}
  \braket{e^{ik'x}}{e^{ikx}}
    =\frac{1}{N}\int_{-\pi}^{\pi}e^{-ik'x}e^{ikx}\dd x .
\end{equation}
La integral se calcula directamente,
\begin{equation}
  \label{eq:integral_ortogonalidad}
  \int_{-\pi}^{\pi}e^{i(k-k')x}\dd x=
  \begin{cases}
    2\pi, & k=k',\\
    0, & k\neq k',
  \end{cases}
  =2\pi\de_{kk'},
\end{equation}
de modo que la normalizaci\'on adecuada es $N=2\pi$ y
\begin{equation}
  \label{eq:ortonormalidad_fourier}
  \braket{e^{ik'x}}{e^{ikx}}
    =\frac{1}{2\pi}\int_{-\pi}^{\pi}e^{-ik'x}e^{ikx}\dd x=\de_{kk'} .
\end{equation}
Es decir, respecto del producto interno
\begin{equation}
  \label{eq:pi_fourier}
  \braket{f}{g}=\frac{1}{2\pi}\int_{-\pi}^{\pi}f^{*}(x)g(x)\dd x,
\end{equation}
la familia $\set{e^{ikx}}_{k=-\infty}^{\infty}$ es una \textbf{base
ortonormal}.

\subsection{Los coeficientes}
\label{subsec:coeficientes_fourier}

La ortonormalidad permite extraer los coeficientes proyectando, exactamente
como en \eqref{eq:contraccion_kronecker}:
\begin{equation}
  \label{eq:extraccion_coeficientes}
  \begin{aligned}
    \braket{e^{ik'x}}{f(x)}
      &=\frac{1}{2\pi}\int_{-\pi}^{\pi}e^{-ik'x}
        \sum_{k=-\infty}^{\infty}C_ke^{ikx}\dd x\\
      &=\sum_{k=-\infty}^{\infty}C_k\,
        \frac{1}{2\pi}\int_{-\pi}^{\pi}e^{-ik'x}e^{ikx}\dd x\\
      &=\sum_{k=-\infty}^{\infty}C_k\de_{kk'}=C_{k'} ,
  \end{aligned}
\end{equation}
es decir
\begin{equation}
  \label{eq:coeficientes_fourier}
  C_k=\braket{e^{ikx}}{f(x)}
    =\frac{1}{2\pi}\int_{-\pi}^{\pi}e^{-ikx}f(x)\dd x .
\end{equation}

\subsection{Forma real}
\label{subsec:fourier_real}

Agrupando en \eqref{eq:serie_fourier} los t\'erminos $\pm k$ y usando la
f\'ormula de Euler \eqref{eq:euler},
\begin{equation}
  \label{eq:fourier_real_desarrollo}
  \small
  \begin{aligned}
    f(x)&=\dots+C_{-1}e^{-ix}+C_0+C_1e^{ix}+\dots\\
        &=C_0+\dots+C_{-1}\qty{\cos x-i\sin x}\\
        &\qquad\quad+C_{1}\qty{\cos x+i\sin x}+\dots\\
        &=C_0+\dots+(C_{1}+C_{-1})\cos x\\
        &\qquad\quad+i(C_{1}-C_{-1})\sin x+\dots
  \end{aligned}
\end{equation}
Definiendo $a_k=C_k+C_{-k}$ y $b_k=i(C_k-C_{-k})$ se obtiene la
\textbf{forma real} de la serie,
\begin{equation}
  \label{eq:fourier_real}
  f(x)=\frac{a_0}{2}+\sum_{k=1}^{\infty}a_k\cos(kx)
    +\sum_{k=1}^{\infty}b_k\sin(kx).
\end{equation}

Las familias $\set{\cos(kx)}_{k=0}^{\infty}$ y
$\set{\sin(kx)}_{k=1}^{\infty}$ son tambi\'en ortonormales, ahora respecto de
la normalizaci\'on $N=\pi$. En efecto, usando la identidad producto-suma,
\begin{equation}
  \label{eq:ortonormalidad_senos}
  \small
  \begin{aligned}
    \braket{\sin(k'x)}{\sin(kx)}
      &=\frac{1}{\pi}\int_{-\pi}^{\pi}\sin(k'x)\sin(kx)\dd x\\
      &=\frac{1}{2\pi}\int_{-\pi}^{\pi}
        \big(\cos((k-k')x)-\cos((k+k')x)\big)\dd x\\
      &=\de_{kk'},
  \end{aligned}
\end{equation}
y an\'alogamente
\begin{equation}
  \label{eq:ortonormalidad_cosenos}
  \braket{\cos(k'x)}{\cos(kx)}=\de_{kk'},
  \qquad
  \braket{\cos(k'x)}{\sin(kx)}=0
\end{equation}
para todos $k,k'\ge 1$. Por lo tanto, proyectando igual que en
\eqref{eq:extraccion_coeficientes},
\begin{equation}
  \label{eq:coeficientes_reales}
  \begin{aligned}
    a_k&=\braket{\cos(kx)}{f(x)}
        =\frac{1}{\pi}\int_{-\pi}^{\pi}f(x)\cos(kx)\dd x,\\
    b_k&=\braket{\sin(kx)}{f(x)}
        =\frac{1}{\pi}\int_{-\pi}^{\pi}f(x)\sin(kx)\dd x,\\
    \frac{a_0}{2}&=\frac{1}{2\pi}\int_{-\pi}^{\pi}f(x)\dd x ,
  \end{aligned}
\end{equation}
donde el \'ultimo t\'ermino es simplemente el valor medio de $f$ sobre el
periodo.
